\documentclass[aspectratio=169,11pt]{beamer}

\usetheme{Madrid}
\usecolortheme{wolverine}

\usepackage{fontspec}
\usepackage{booktabs}
\usepackage{array}
\usepackage{hyperref}
\usepackage{mihajlo-elfak-logos}

\setsansfont{Latin Modern Sans}
\setmonofont{Latin Modern Mono}
\hypersetup{colorlinks=true, urlcolor=blue}
\setbeamertemplate{navigation symbols}{}

\title{Projekat I}
\subtitle[AIS analitika nad HDFS + Spark]{AIS analitika nad HDFS + Spark}
\author[Mihajlo Madić 2119]{Mihajlo Madić 2119}
\institute[Big Data Systems]{Big Data Systems}
\date{\small{\textit{Jun 2026, Niš}}}

\newcommand{\code}[1]{\texttt{\small #1}}

\begin{document}

\begin{frame}
  \titlepage
\end{frame}

\begin{frame}{Dataset: Piraeus AIS}
  \begin{itemize}
    \item Korišćen je javni dataset \textit{The Piraeus AIS Dataset for Large-scale Maritime Data Analytics}, objavljen u okviru Zenodo repozitorijuma.
    \item Dataset sadrži AIS\footnote{Automatic Identification System.} zapise kreirane tokom aktivnosti brodova u okviru luke Pirej u Grčkoj.
    \item AIS zapisi opisuju kretanje brodova kroz atribute kao što su vreme slanja (\code{timestamp}), identifikator broda (\code{vessel\_id}), geografska pozicija (\code{lon}, \code{lat}), brzina (\code{speed}), kurs\footnote{Kurs (\code{course}) opisuje stvarni smer kretanja broda, dok pravac (\code{heading}) opisuje smer njegovog pramca.} (\code{course}) i pravac (\code{heading}).
    \item Originalni dataset je dovoljno veliki za Big Data scenario; za ovaj prototip korišćen je kontrolisani podskup radi bržeg razvoja i lakšeg i kraćeg testiranja, i generisanja benchmark-ova.
    \item Izvor: \url{https://zenodo.org/records/5792100}
  \end{itemize}
\end{frame}

\begin{frame}{Zadatak}
  Sistem:
  \begin{enumerate}
    \item priprema podataka i upload na HDFS
    \item Spark aplikacija sa CLI upitima (\code{count}, \code{show}, \code{stats})
    \item benchmark poređenje režima izvršavanja
  \end{enumerate}

  \vspace{0.4cm}
  Rezultati:
  \begin{itemize}
    \item reproducibilna infrastruktura kroz Docker Compose
    \item merenja performansi
    \item izveštaji sa rezultatima
  \end{itemize}
\end{frame}

\begin{frame}{Obim i podaci}
  Korišćen podskup:
  \begin{itemize}
    \item raw podaci: \code{\textasciitilde 971.54 MiB} (jul + decembar 2018 CSV)
    \item preprocesirani podaci: \code{\textasciitilde 191.60 MiB} (Parquet)
  \end{itemize}

  \vspace{0.35cm}
  Relevantne HDFS putanje:
  \begin{itemize}
    \item \code{/bds/proj1/raw}
    \item \code{/bds/proj1/processed/clean}
    \item \code{/bds/proj1/processed/quarantine}
    \item \code{/bds/proj1/processed/quality\_report}
  \end{itemize}
\end{frame}

\begin{frame}{Arhitektura}
  3-slojni model:
  \begin{enumerate}
    \item HDFS sloj: skladište (\code{namenode}, \code{datanode})
    \item Spark sloj: obrada (\code{spark-master}, \code{spark-worker-*})
    \item App sloj: job submission i CLI (\code{app})
  \end{enumerate}

  \vspace{0.35cm}
  \begin{block}{Tok podataka}
    Raw CSV $\rightarrow$ HDFS \code{/raw} $\rightarrow$ preprocess job $\rightarrow$ HDFS \code{/processed/*}
  \end{block}

  \begin{itemize}
    \item analytics CLI koristi \code{count}, \code{show}, \code{stats}
    \item svi benchmark testovi koriste isti izvor podataka (HDFS)
  \end{itemize}
\end{frame}

\begin{frame}{Docker Compose}
  Servisi:
  \begin{itemize}
    \item \code{namenode}, \code{datanode} \Rightarrow \space storage
    \item \code{spark-master}, \code{spark-worker-1}, \code{spark-worker-2} \Rightarrow \space compute
    \item \code{app} \Rightarrow \space "bussiness" logic
  \end{itemize}

  \vspace{0.35cm}
  Princip:
  \begin{itemize}
    \item isti HDFS izvor za oba benchmark moda
    \item \code{app} container submituje iste komande za \code{local[*]} i \code{spark://...}
  \end{itemize}
\end{frame}

\begin{frame}[fragile]{Prva faza - upload podataka na HDFS}
  Checkpoint:
  \begin{center}
	\code{automation/phase1\_hdfs\_upload.sh}
  \end{center}

  Ključne komande:
  \vspace{0.15cm}
  \scriptsize
  \begin{tabular}{p{0.46\textwidth}p{0.46\textwidth}}
    \toprule
    Komanda & Šta radi \\
    \midrule
    \code{docker compose exec namenode} &
    Izvršava HDFS komande unutar \code{namenode} kontejnera, tj. tamo gde je dostupan Hadoop CLI. \\
    \code{hdfs dfs -mkdir -p /bds/proj1/raw} &
    Kreira raw direktorijum u HDFS-u; opcija \code{-p} sprečava grešku ako putanja već postoji. \\
    \code{hdfs dfs -put -f /input/.../jul2018.csv /bds/proj1/raw/jul2018.csv} &
    Uploaduje CSV za jul 2018. iz read-only input mount-a u HDFS raw zonu; \code{-f} omogućava overwrite. \\
    \code{hdfs dfs -put -f /input/.../dec2018.csv /bds/proj1/raw/dec2018.csv} &
    Uploaduje CSV za decembar 2018. na istu HDFS lokaciju. \\
    \code{hdfs dfs -count -h /bds/proj1/raw} &
    Verifikuje upload prikazom broja direktorijuma, fajlova i ukupne veličine u formatu čitljivom za čoveka. \\
    \bottomrule
  \end{tabular}
  \normalsize
\end{frame}

\begin{frame}[fragile]{Druga faza - preprocesiranje podataka (kroz Spark job)}
  Checkpoint:
\begin{center}
	\code{automation/phase2\_preprocessing.sh}
\end{center}
  
Job:
  \begin{verbatim}
docker compose exec app /spark/bin/spark-submit \
  --master spark://spark-master:7077 \
  /workspace/scripts/preprocess_ais_hdfs.py \
  --input 'hdfs://namenode:9000/bds/proj1/raw/*.csv' \
  --output-base hdfs://namenode:9000/bds/proj1/processed
  \end{verbatim}
\end{frame}

\begin{frame}[fragile]{Preprocessing logika}
  Ključne transformacije:
  \begin{verbatim}
missing_required = timestamp/vessel_id/lon/lat missing
invalid_geo      = lon/lat out of range
invalid_speed    = speed < 0
invalid_heading  = heading not in [0, 360]
invalid_course   = course not in [0, 360]

clean       <- no hard_fail
quarantine  <- hard_fail rows

clean:       parquet partitioned by year, month
quarantine:  parquet partitioned by reason
quality_report: summary + reason_counts
  \end{verbatim}
\end{frame}

\begin{frame}[fragile]{Analytics CLI}
  Komande:
  \begin{itemize}
    \item \code{count}
    \item \code{show}
    \item \code{stats} (min/max/avg/stddev + count)
  \end{itemize}

  Primeri:
  \begin{verbatim}
ais_analytics.py count --year 2018 --month 12

ais_analytics.py show --year 2018 --month 7 --min-speed 15 \
  --sort-by timestamp --sort-desc --limit 20

ais_analytics.py stats --year 2018 --group-by month --metrics speed
  \end{verbatim}
\end{frame}

\begin{frame}{Benchmark metodologija}
  Pravila za fer poređenje:
  \begin{itemize}
    \item isti dataset: \code{/processed/clean}
    \item ista komanda
    \item režimi: \code{--master local[*]} i \code{--master spark://spark-master:7077}
  \end{itemize}

  \vspace{0.35cm}
  Kampanja (jedan kompletan benchmark):
  \begin{itemize}
    \item 5 iteracija
    \item svaka iteracija: 1 izvršenje po modu (klaster/standalone) + zaseban izveštaj
    \item zatim finalni agregat (5x2)
  \end{itemize}
\end{frame}

\begin{frame}[fragile]{Benchmark alati i artefakti}
  Runner:
  \begin{verbatim}
python3 scripts/benchmark_ais_analytics.py \
  --benchmark-name <name> \
  --modes both \
  --repeats 1 \
  --query-override "count ... | show ... | stats ..."
  \end{verbatim}

  Aggregator:
  \begin{verbatim}
python3 scripts/aggregate_benchmark_runs.py \
  --runs-dir runs/campaign_<id>/runs \
  --output-root runs/campaign_<id>/reports
  \end{verbatim}
\end{frame}

\begin{frame}{Finalni sistem}
  Checkpoint:
  \begin{center}
    \code{automation/phase5\_full\_benchmarks.sh}
  \end{center}

  \vspace{0.35cm}
  Šta radi:
  \begin{enumerate}
    \item podiže storage + compute + app sloj
    \item puni HDFS i izvršava preprocessing
    \item pokreće 4 benchmark kampanje
    \item za svaku kampanju pravi 5 iteracionih reporta + 1 finalni agregat
  \end{enumerate}
\end{frame}

\begin{frame}[fragile]{Struktura rezultata}
  \scriptsize
  \begin{verbatim}
runs/campaign_<id>/
  runs/
    standalone/<timestamp>/
      meta.json
      results.json
      summary.json
      run_report.md
    cluster/<timestamp>/
      ...
  reports/
    <timestamp>/aggregate.md
    <timestamp>/aggregate.csv
    <timestamp>/aggregate.json
  \end{verbatim}
  \normalsize

  \code{meta.json} čuva raw i processed size snapshot.
\end{frame}

\begin{frame}[fragile]{Rezultati 1GB: Stats (speed)}
  Kampanja: \code{campaign\_20260226T101745Z\_stats\_speed}

  \scriptsize
  \begin{verbatim}
stats --year 2018 --group-by month --metrics speed \
  --order-by month --limit 12
  \end{verbatim}
  \normalsize

  \begin{tabular}{lrrrr}
    \toprule
    mode & n & mean [s] & min [s] & max [s] \\
    \midrule
    standalone & 5 & 11.049 & 10.350 & 11.781 \\
    cluster & 5 & 15.444 & 15.023 & 16.202 \\
    \bottomrule
  \end{tabular}
\end{frame}

\begin{frame}[fragile]{Rezultati 1GB: Stats (heading, course)}
  Kampanja: \code{campaign\_20260226T102011Z\_stats\_heading\_course}

  \scriptsize
  \begin{verbatim}
stats --year 2018 --group-by month --metrics heading,course \
  --order-by month --limit 12
  \end{verbatim}
  \normalsize

  \begin{tabular}{lrrrr}
    \toprule
    mode & n & mean [s] & min [s] & max [s] \\
    \midrule
    standalone & 5 & 11.132 & 10.760 & 11.825 \\
    cluster & 5 & 16.069 & 15.484 & 16.734 \\
    \bottomrule
  \end{tabular}

\end{frame}

\begin{frame}[fragile]{Rezultati 1GB: Count}
  Kampanja: \code{campaign\_20260226T102241Z\_count}

  \begin{verbatim}
count --year 2018 --month 12
  \end{verbatim}

  \begin{tabular}{lrrrr}
    \toprule
    mode & n & mean [s] & min [s] & max [s] \\
    \midrule
    standalone & 5 & 9.486 & 9.465 & 9.508 \\
    cluster & 5 & 14.065 & 13.657 & 14.188 \\
    \bottomrule
  \end{tabular}
\end{frame}

\begin{frame}[fragile]{Rezultati 1GB: Show}
  Kampanja: \code{campaign\_20260226T102453Z\_show}

  \scriptsize
  \begin{verbatim}
show --year 2018 --month 7 --min-speed 15 \
  --sort-by timestamp --sort-desc --limit 20
  \end{verbatim}
  \normalsize

  \begin{tabular}{lrrrr}
    \toprule
    mode & n & mean [s] & min [s] & max [s] \\
    \midrule
    standalone & 5 & 10.490 & 10.445 & 10.513 \\
    cluster & 5 & 14.309 & 14.167 & 14.697 \\
    \bottomrule
  \end{tabular}
\end{frame}

\begin{frame}{Poređenje: 1GB benchmark}
  Cluster/standalone odnos za \textbf{manji} skup:

  \vspace{0.25cm}
  \scriptsize
  \begin{tabular}{lrrrr}
    \toprule
    benchmark & standalone mean [s] & cluster mean [s] & ratio & delta [s] \\
    \midrule
    stats speed & 11.049 & 15.444 & 1.397x & +4.395 \\
    stats heading, course & 11.132 & 16.069 & 1.443x & +4.937 \\
    count dec2018 & 9.486 & 14.065 & 1.483x & +4.579 \\
    show fast jul2018 & 10.490 & 14.309 & 1.364x & +3.819 \\
    \bottomrule
  \end{tabular}
  \normalsize

  \vspace{0.3cm}
  Raw: \code{971.54 MiB}, processed parquet: \code{191.60 MiB}, clean rows: \code{8.30M}
\end{frame}

\begin{frame}[fragile]{Rezultati 5GB: Stats (speed)}
  Kampanja: \code{campaign\_20260627T140125Z\_5gb\_stats\_speed}

  \scriptsize
  \begin{verbatim}
stats --year 2018 --group-by month --metrics speed \
  --order-by month --limit 12
  \end{verbatim}
  \normalsize

  \begin{tabular}{lrrrr}
    \toprule
    mode & n & mean [s] & min [s] & max [s] \\
    \midrule
    standalone & 5 & 12.882 & 12.379 & 13.497 \\
    cluster & 5 & 17.692 & 17.021 & 18.185 \\
    \bottomrule
  \end{tabular}

\end{frame}

\begin{frame}[fragile]{Rezultati 5GB: Stats (heading, course)}
  Kampanja: \code{campaign\_20260627T140416Z\_5gb\_stats\_heading\_course}

  \scriptsize
  \begin{verbatim}
stats --year 2018 --group-by month --metrics heading,course \
  --order-by month --limit 12
  \end{verbatim}
  \normalsize

  \begin{tabular}{lrrrr}
    \toprule
    mode & n & mean [s] & min [s] & max [s] \\
    \midrule
    standalone & 5 & 14.551 & 14.493 & 14.612 \\
    cluster & 5 & 19.803 & 19.196 & 20.193 \\
    \bottomrule
  \end{tabular}

\end{frame}

\begin{frame}[fragile]{Rezultati 5GB: Count}
  Kampanja: \code{campaign\_20260627T140724Z\_5gb\_count}

  \begin{verbatim}
count --year 2018 --month 12
  \end{verbatim}

  \begin{tabular}{lrrrr}
    \toprule
    mode & n & mean [s] & min [s] & max [s] \\
    \midrule
    standalone & 5 & 11.010 & 10.967 & 11.047 \\
    cluster & 5 & 15.795 & 15.666 & 16.215 \\
    \bottomrule
  \end{tabular}
\end{frame}

\begin{frame}[fragile]{Rezultati 5GB: Show}
  Kampanja: \code{campaign\_20260627T140955Z\_5gb\_show}

  \scriptsize
  \begin{verbatim}
show --year 2018 --month 7 --min-speed 15 \
  --sort-by timestamp --sort-desc --limit 20
  \end{verbatim}
  \normalsize

  \begin{tabular}{lrrrr}
    \toprule
    mode & n & mean [s] & min [s] & max [s] \\
    \midrule
    standalone & 5 & 12.035 & 12.004 & 12.052 \\
    cluster & 5 & 16.529 & 16.153 & 16.815 \\
    \bottomrule
  \end{tabular}
\end{frame}

\begin{frame}{Rezultati: \textasciitilde 5GB benchmark}
  \begin{itemize}
    \item Ulazni skup podataka je proširen dodatnim podskupovima grupisanim po mesecima.
    \item Cluster/standalone odnos za \textbf{veći} skup:
  \end{itemize}

  \vspace{0.5cm}
  \scriptsize
  \begin{tabular}{lrrrr}
    \toprule
    benchmark & standalone mean [s] & cluster mean [s] & ratio & delta [s] \\
    \midrule
    stats speed & 12.882 & 17.692 & 1.373x & +4.810 \\
    stats heading, course & 14.551 & 19.803 & 1.361x & +5.252 \\
    count dec2018 & 11.010 & 15.795 & 1.435x & +4.785 \\
    show fast jul2018 & 12.035 & 16.529 & 1.373x & +4.494 \\
    \bottomrule
  \end{tabular}
  \normalsize

  \vspace{0.3cm}
  Raw: \code{4.60 GiB}, processed parquet: \code{812.68 MiB}, clean rows: \code{40.29M}
\end{frame}

\begin{frame}{Skaliranje: 1GB prema 5GB}
  Poređenje srednjih vremena izvršavanja:

  \vspace{0.15cm}
  \scriptsize
  \begin{tabular}{lrrrr}
    \toprule
    benchmark & 1GB standalone & 1GB cluster & 5GB standalone & 5GB cluster \\
    \midrule
    stats speed & 11.049 & 15.444 & 12.882 & 17.692 \\
    stats heading, course & 11.132 & 16.069 & 14.551 & 19.803 \\
    count dec2018 & 9.486 & 14.065 & 11.010 & 15.795 \\
    show fast jul2018 & 10.490 & 14.309 & 12.035 & 16.529 \\
    \bottomrule
  \end{tabular}
  \normalsize

  \vspace{0.3cm}
  Odnos veličina:
  \begin{itemize}
    \item raw: \code{971.54 MiB} prema \code{4.60 GiB} (\textasciitilde 4.85x)
    \item parquet: \code{191.60 MiB} prema \code{812.68 MiB} (\textasciitilde 4.24x)
  \end{itemize}
\end{frame}

\begin{frame}{Analiza: manji benchmark}
  Zaključci za manji obim (\textasciitilde 1GB raw, \textasciitilde 192MB parquet):
  \begin{itemize}
    \item \code{local[*]} je brži u svim testovima.
    \item Upiti su kratki i latency-bound: skeniranje parquet-a, filtriranje i male agregacije.
    \item Cluster dodaje overhead: scheduling, driver/executor komunikaciju i serijalizaciju (Docker mreža).
    \item \code{count} i \code{show} su stabilniji zbog jednostavnosti; \code{stats} ima više varijanse zbog grupisanja i agregacija.
  \end{itemize}
\end{frame}

\begin{frame}{Analiza: veći benchmark}
  Rezultat je očekivan u odnosu na manji benchmark:
  \begin{itemize}
    \item Cluster je i dalje sporiji, a delta ostaje oko \code{+4.5s} do \code{+5.3s}.
    \item Odnos vremena se nije bitno promenio: \code{1.36x} do \code{1.44x}.
    \item Veći ulaz povećava vreme, \textasciitilde linearno sa veličinom ulaznog skupa (parquet format).
    \item Fiksni Spark overhead još uvek dominira nad dobicima zbog paralelizacije.
  \end{itemize}
\end{frame}

\begin{frame}{Zašto cluster nema bolje performanse?}
  Infrastruktura:
  \begin{itemize}
    \item Cluster radi na istom fizičkom hostu, sa ograničenim brojem worker-a.
    \item HDFS ima jedan datanode, pa nema prave distribuirane data locality prednosti.
    \item Worker-i i Spark master komuniciraju preko Docker mreže.
    \item Testirani upiti nisu dovoljno teški da amortizuju cenu distribuiranog izvršavanja.
  \end{itemize}

  \vspace{0.3cm}
  Najveći efekat skaliranja se video u preprocessing fazi: čitanje \code{4.60 GiB} raw CSV-a,
  validacija \code{40.29M} redova i upis u storage sistem (HDFS).
\end{frame}

\begin{frame}{Ograničenja i sledeći koraci}
  Ograničenja eksperimenta:
  \begin{itemize}
    \item dva obima podataka, ali jedan host i jedan datanode
    \item meren je latency pojedinačnih batch upita, ne throughput pod konkurentnim opterećenjem
    \item workload koristi relativno jednostavne filtre, sortiranje i agregacije
  \end{itemize}

  \vspace{0.3cm}
  Za očekivanu cluster prednost može se testirati:
  \begin{itemize}
    \item više datanode/worker čvorova
    \item veći procesirani skup i teže operacije
    \item Promena Spark cluster konfiguracije: executor memorija i broj executor-a
  \end{itemize}
\end{frame}

\begin{frame}[fragile]{Ponavljanje eksperimenta -- referenca}

  \begin{verbatim}
# HDFS + preprocess baseline
bash automation/phase2_preprocessing.sh

# jedna benchmark kampanja
python3 scripts/benchmark_ais_analytics.py ...
python3 scripts/aggregate_benchmark_runs.py ...

# 5GB benchmark sa cleanup-om
bash automation/phase5_5gb_benchmarks.sh
  \end{verbatim}

  Sve kampanje i izveštaji su timestamped i append-only.
\end{frame}

\begin{frame}
  \centering
  {\Huge Pitanja?}

  \vspace{0.8cm}
  {\Large Kraj}
\end{frame}

\end{document}
